\chapter{Poisson Brackets, Generators, and Angular Momentum}

\lectureoverview{The Poisson bracket is more than a compact way to rewrite
Hamilton's equations.  It is an algebraic operation that generates
transformations.  Bracketing with the Hamiltonian generates time evolution;
bracketing with momentum generates translations; bracketing with angular
momentum generates rotations.  In this language, symmetry and conservation
become the two readings of one vanishing bracket.  We shall develop that idea
from the harmonic oscillator and then use the angular-momentum algebra to
obtain gyroscopic precession without introducing Euler angles.}

\section{A Function Along a Hamiltonian Trajectory}

Let
\[
  F=F(q_1,\ldots,q_n,p_1,\ldots,p_n)
\]
be a sufficiently differentiable function on phase space.  It can be viewed
in two ways.  First, it assigns a number to every point of phase space.
Second, once a Hamiltonian trajectory has been chosen, it becomes a function
of time because the point \((q(t),p(t))\) moves.  Along that trajectory,
\begin{equation}
  \frac{dF}{dt}
  =
  \sum_i\left(
    \frac{\partial F}{\partial q_i}\dot q_i
    +\frac{\partial F}{\partial p_i}\dot p_i
  \right).
  \label{eq:l8-chain-rule}
\end{equation}
Hamilton's equations,
\begin{equation}
  \dot q_i=\frac{\partial H}{\partial p_i},
  \qquad
  \dot p_i=-\frac{\partial H}{\partial q_i},
  \label{eq:l8-hamilton-equations}
\end{equation}
turn this into
\begin{equation}
  \frac{dF}{dt}
  =
  \sum_i\left(
    \frac{\partial F}{\partial q_i}
    \frac{\partial H}{\partial p_i}
    -
    \frac{\partial F}{\partial p_i}
    \frac{\partial H}{\partial q_i}
  \right).
  \label{eq:l8-fdot-expanded}
\end{equation}
The combination on the right occurs throughout classical mechanics, even when
neither of its entries is the Hamiltonian.  This motivates the definition
\begin{equation}
  \boxed{
  \{F,G\}
  =
  \sum_i\left(
    \frac{\partial F}{\partial q_i}
    \frac{\partial G}{\partial p_i}
    -
    \frac{\partial F}{\partial p_i}
    \frac{\partial G}{\partial q_i}
  \right).}
  \label{eq:l8-poisson-definition}
\end{equation}
Equation~\eqref{eq:l8-fdot-expanded} is therefore
\begin{equation}
  \boxed{\dot F=\{F,H\}.}
  \label{eq:l8-time-generator}
\end{equation}

\begin{classroomqa}
  \audiencequestion{If the second entry is changed from \(H\) to an arbitrary
  function \(G\), is the bracket still the time derivative of \(F\)?
  \lecturetimestamp{00:05:45--00:06:15}}
  \lecturerresponse{No.  The audience catches an important point while the
  notation is being generalized.  Equation~\eqref{eq:l8-poisson-definition}
  defines \(\{F,G\}\) for any two phase-space functions, but
  \(\dot F=\{F,G\}\) holds only when \(G\) is the Hamiltonian.  A general
  \(G\) will generate a different infinitesimal transformation.}
\end{classroomqa}

It is useful to write a small change in time as
\begin{equation}
  \delta_t F=\epsilon\,\{F,H\}.
  \label{eq:l8-small-time-step}
\end{equation}
This form anticipates what comes later.  Time evolution is one member of a
larger family of transformations generated by Poisson brackets.

\section{The Algebra of Brackets}

The bracket can now be used as a calculus of its own.  Its rules follow from
ordinary differentiation:
\begin{align}
  \{A,B\}&=-\{B,A\}, \label{eq:l8-antisymmetry}\\
  \{A+B,C\}&=\{A,C\}+\{B,C\}, \label{eq:l8-additivity}\\
  \{\lambda A,B\}&=\lambda\{A,B\},
    \quad \lambda\ \text{constant}, \label{eq:l8-scalar-rule}\\
  \{AB,C\}&=A\{B,C\}+\{A,C\}B. \label{eq:l8-leibniz}
\end{align}
These four lines express antisymmetry, additivity, linearity under multiplication
by a constant, and the Leibniz rule, respectively.
Classical phase-space functions multiply commutatively, so the placement of
the ordinary factors \(A\) and \(B\) in the last line is immaterial.  That
will no longer be true for operators in quantum mechanics.  Antisymmetry also
implies
\begin{equation}
  \{A,A\}=0.
  \label{eq:l8-self-bracket}
\end{equation}

For canonical coordinates the elementary brackets are
\begin{equation}
  \boxed{
  \{q_i,q_j\}=0,\qquad
  \{p_i,p_j\}=0,\qquad
  \{q_i,p_j\}=\delta_{ij}.}
  \label{eq:l8-canonical-brackets}
\end{equation}
Any function only of the \(q_i\) has zero bracket with any other function only
of the \(q_i\); the analogous statement holds for functions only of the
\(p_i\).  A nonzero bracket requires a canonical coordinate to meet its
conjugate momentum.  The Kronecker delta in
Eq.~\eqref{eq:l8-canonical-brackets} records that only matching indices meet.

\subsection{The harmonic oscillator as a test}

Consider
\begin{equation}
  H=\frac{p^2}{2m}+\frac{kq^2}{2}.
  \label{eq:l8-oscillator-hamiltonian}
\end{equation}
We now deliberately use only the bracket rules and
Eq.~\eqref{eq:l8-time-generator}.  First,
\begin{align}
  \dot q
  &=\{q,H\}\\
  &=\frac{1}{2m}\{q,p^2\}\\
  &=\frac{1}{2m}\left(p\{q,p\}+\{q,p\}p\right)\\
  &=\frac{p}{m}.
  \label{eq:l8-qdot-oscillator}
\end{align}
The potential term disappears because \(\{q,q^2\}=0\).  Similarly,
\begin{align}
  \dot p
  &=\{p,H\}\\
  &=\frac{k}{2}\{p,q^2\}\\
  &=\frac{k}{2}\left(q\{p,q\}+\{p,q\}q\right)\\
  &=-kq.
  \label{eq:l8-pdot-oscillator}
\end{align}
The minus sign is precisely the antisymmetry
\(\{p,q\}=-\{q,p\}=-1\).  Combining the two first-order equations gives
\begin{equation}
  m\ddot q=-kq.
  \label{eq:l8-hooke-equation}
\end{equation}

\begin{figure}[htbp]
  \centering
  \includegraphics[width=0.82\textwidth]{lecture_08_figure_01.png}
  \caption{The oscillator Hamiltonian and the two bracket calculations
  yielding \(\dot q=p/m\) and \(\dot p=-kq\).}
  \lectureframe{00:23:20}
\end{figure}

During the calculation the class looks for a verb meaning
\emph{take a Poisson bracket} and settles, provisionally, on
\emph{Poissonate}.  The joke marks a real change of viewpoint: once the
algebra is known, one can work with it directly instead of re-deriving every
result from the definition.

\begin{classroomqa}
  \audiencequestion{Is this mathematical structure a group, or an algebra?
  Does it have inverses? \lecturetimestamp{00:23:48--00:24:16}}
  \lecturerresponse{It is not a group: the bracket operation does not provide
  the required identity and inverses.  The phase-space functions form a
  Poisson algebra: they have ordinary commutative multiplication and addition,
  together with an antisymmetric bracket obeying the Leibniz rule.  A further
  identity, the Jacobi identity, completes the usual abstract definition,
  although it is not developed in this lecture.}
\end{classroomqa}

\begin{classroomqa}
  \audiencequestion{What is the connection between Poisson brackets and
  quantum commutators? \lecturetimestamp{00:24:18--00:25:00}}
  \lecturerresponse{The lecture gives only a preview.  In the passage to
  quantum mechanics, the classical bracket is replaced schematically by
  \[
    \{A,B\}_{\mathrm{PB}}
    \longleftrightarrow
    \frac{1}{i\hbar}[\widehat A,\widehat B].
  \]
  This correspondence requires qualifications about operator ordering and
  quantization; it is not being used as a theorem here.}
\end{classroomqa}

For the oscillator, the bracket has merely reconstructed Hamilton's equations.
Its real economy appears when the generator is a conserved quantity other than
\(H\).

\section{Angular Momentum Recovered from Rotational Symmetry}

Recall the infinitesimal form of Noether's construction used earlier in the
course.  If a continuous symmetry changes the coordinates by
\begin{equation}
  \delta q_i=\epsilon F_i(q),
  \label{eq:l8-noether-variation}
\end{equation}
then, in the setting considered here, its conserved quantity is
\begin{equation}
  G=\sum_i p_iF_i(q).
  \label{eq:l8-noether-charge}
\end{equation}
The overall factor \(\epsilon\) may be omitted: multiplying a conserved
quantity by a fixed number does not change the conservation law.

Take one particle and rotate it infinitesimally about the \(z\)-axis.  With the
orientation used in the lecture,
\begin{equation}
  \delta x=-\epsilon y,\qquad
  \delta y=\epsilon x,\qquad
  \delta z=0.
  \label{eq:l8-z-rotation}
\end{equation}
Equation~\eqref{eq:l8-noether-charge} gives
\begin{equation}
  L_z=-p_xy+p_yx=xp_y-yp_x.
  \label{eq:l8-lz}
\end{equation}
Cycling \(x\), \(y\), and \(z\),
\begin{equation}
  \boxed{\begin{aligned}
    L_x&=yp_z-zp_y,\\
    L_y&=zp_x-xp_z,\\
    L_z&=xp_y-yp_x.
  \end{aligned}}
  \label{eq:l8-angular-components}
\end{equation}
These are the components of
\begin{equation}
  \boldsymbol L=\boldsymbol r\times\boldsymbol p.
  \label{eq:l8-cross-product}
\end{equation}

\begin{figure}[htbp]
  \centering
  \includegraphics[width=0.82\textwidth]{lecture_08_figure_02.png}
  \caption{The three Cartesian components of angular momentum and their
  compact cross-product form.}
  \lectureframe{00:34:00}
\end{figure}

Rotation about an arbitrary unit vector \(\boldsymbol n\) is generated by
\begin{equation}
  L_{\boldsymbol n}=\boldsymbol n\cdot\boldsymbol L.
  \label{eq:l8-arbitrary-axis}
\end{equation}
For an isolated rotationally invariant system, all three components are
conserved.  If only rotations around one axis remain a symmetry, only the
component along that axis is guaranteed to be conserved.

\section{Angular Momentum Generates Rotations}

Instead of immediately solving a rotational dynamics problem, let us explore
what \(L_z\) does under the bracket.  Using
Eq.~\eqref{eq:l8-lz},
\begin{align}
  \{x,L_z\}
  &=\{x,xp_y-yp_x\}
   =-y\{x,p_x\}
   =-y,\\
  \{y,L_z\}
  &=\{y,xp_y-yp_x\}
   =x\{y,p_y\}
   =x,\\
  \{z,L_z\}&=0.
  \label{eq:l8-position-rotation-brackets}
\end{align}
These are exactly the coefficients in
Eq.~\eqref{eq:l8-z-rotation}.

\begin{figure}[htbp]
  \centering
  \includegraphics[width=0.82\textwidth]{lecture_08_figure_03.png}
  \caption{The coordinate brackets with \(L_z\): the bracket reproduces the
  infinitesimal rotation of \((x,y,z)\).}
  \lectureframe{00:38:41}
\end{figure}

Momentum is also a vector, and the same calculation gives
\begin{equation}
  \begin{aligned}
    \{p_x,L_z\}&=-p_y,\\
    \{p_y,L_z\}&=p_x,\\
    \{p_z,L_z\}&=0.
  \end{aligned}
  \label{eq:l8-momentum-rotation-brackets}
\end{equation}

\begin{figure}[htbp]
  \centering
  \includegraphics[width=0.82\textwidth]{lecture_08_figure_04.png}
  \caption{Position and momentum transform with the same vector pattern under
  the rotation generated by \(L_z\).}
  \lectureframe{00:42:08}
\end{figure}

The general Hamiltonian statement is therefore
\begin{equation}
  \boxed{\delta_G A=\epsilon\{A,G\}.}
  \label{eq:l8-generator-rule}
\end{equation}
Every phase-space function \(G\) defines such an infinitesimal canonical
transformation.  If \(G=L_z\), it is a rotation.  If \(G=H\), it is a small
step in time.

\begin{classroomqa}
  \audiencequestion{The lecture sometimes speaks of rotating the system and
  sometimes of rotating coordinates.  Which convention is used in the bracket
  equations? \lecturetimestamp{00:28:22--00:29:35}}
  \lecturerresponse{Equation~\eqref{eq:l8-z-rotation} fixes the convention:
  it is the infinitesimal active rotation
  \((x,y)\mapsto(x-\epsilon y,y+\epsilon x)\).  A passive rotation of the axes
  carries the opposite sign.  Once the convention is fixed, the bracket
  formulas are unambiguous.}
\end{classroomqa}

\section{Momentum Generates Translations}

A translation in one coordinate has
\begin{equation}
  \delta q=\epsilon.
  \label{eq:l8-q-translation}
\end{equation}
Here \(F(q)=1\) in Eq.~\eqref{eq:l8-noether-variation}, so the corresponding
conserved quantity is \(p\).  The generator claim predicts
\begin{equation}
  \{f(q),p\}=\frac{df}{dq}.
  \label{eq:l8-translation-derivative}
\end{equation}
The lecture proves this first for powers by induction.  The base case is
\(\{q,p\}=1\).  Suppose
\begin{equation}
  \{q^{n-1},p\}=(n-1)q^{n-2}.
  \label{eq:l8-induction-hypothesis}
\end{equation}
Then the product rule gives
\begin{align}
  \{q^n,p\}
  &=\{q\,q^{n-1},p\}\\
  &=q\{q^{n-1},p\}+\{q,p\}q^{n-1}\\
  &=(n-1)q^{n-1}+q^{n-1}\\
  &=nq^{n-1}.
  \label{eq:l8-power-translation}
\end{align}
Linearity extends the result to polynomials and to convergent power-series
representations.  More directly, the definition of the bracket gives
Eq.~\eqref{eq:l8-translation-derivative} for every differentiable \(f(q)\):
\begin{equation}
  \{f(q),p\}
  =
  \frac{\partial f}{\partial q}
  \frac{\partial p}{\partial p}
  =f'(q).
\end{equation}

\begin{figure}[htbp]
  \centering
  \includegraphics[width=0.82\textwidth]{lecture_08_figure_05.png}
  \caption{The translation result after the induction argument:
  \(\{F(q),p\}=dF/dq\).}
  \lectureframe{00:53:18}
\end{figure}

We have now seen three versions of one idea:
\begin{center}
  \begin{tabular}{@{}>{\raggedright\arraybackslash}p{0.37\textwidth}
                      >{\raggedright\arraybackslash}p{0.55\textwidth}@{}}
    \toprule
    Generator and transformation & Infinitesimal action\\
    \midrule
    \(H\), time translation & \(\delta_tA=\epsilon\{A,H\}\)\\
    \(p_i\), spatial translation & \(\delta A=\epsilon\{A,p_i\}\)\\
    \(L_i\), rotation & \(\delta A=\epsilon\{A,L_i\}\)\\
    \bottomrule
  \end{tabular}
\end{center}
The conserved quantity associated with a continuous symmetry is also the
generator of that symmetry.

\section{One Bracket, Two Readings}

Let \(G(q,p)\) be a proposed generator.  If it is conserved and has no explicit
time dependence, then
\begin{equation}
  \dot G=\{G,H\}=0.
  \label{eq:l8-g-conserved}
\end{equation}
By antisymmetry,
\begin{equation}
  \{H,G\}=0.
  \label{eq:l8-h-invariant}
\end{equation}
The equations are mathematically equivalent, but their physical readings are
different:
\begin{itemize}
  \item \(\{G,H\}=0\) says that \(G\) does not change under time evolution
  generated by \(H\).
  \item \(\{H,G\}=0\) says that \(H\) does not change under the transformation
  generated by \(G\).
\end{itemize}
The first is a conservation law; the second is a symmetry of the Hamiltonian.
Noether's connection has been compressed into one zero.

\begin{figure}[htbp]
  \centering
  \includegraphics[width=0.82\textwidth]{lecture_08_figure_06.png}
  \caption{The two orders of the same vanishing bracket: conservation of
  \(G\) and invariance of \(H\) under the \(G\)-generated transformation.}
  \lectureframe{01:03:05}
\end{figure}

\begin{classroomqa}
  \audiencequestion{If a conserved quantity is known, can the associated
  symmetry always be recovered? \lecturetimestamp{01:01:42--01:02:15}}
  \lecturerresponse{Within this Hamiltonian setting, yes: bracket every
  phase-space variable with \(G\).  The collection
  \(\delta_GA=\epsilon\{A,G\}\) tells how the variables transform.  For
  example, the brackets with \(L_z\) reveal a rotation even if that
  interpretation was not known in advance.  Global and boundary subtleties
  can arise in more advanced settings, but the local infinitesimal statement
  is the one used here.}
\end{classroomqa}

\begin{classroomqa}
  \audiencequestion{Do those bracket equations tell us the physical motion
  responsible for the conservation law?
  \lecturetimestamp{00:56:01--00:58:13}}
  \lecturerresponse{No.  They identify the symmetry transformation, not the
  actual trajectory.  Conservation of \(L_z\) indicates rotational symmetry;
  it does not say that the particle moves in a circle.  Actual time evolution
  is generated by \(H\).  A bracket with another generator describes another
  direction of transformation in phase space.}
\end{classroomqa}

\begin{classroomqa}
  \audiencequestion{Is there a general transformation statement that does not
  already assume conservation? \lecturetimestamp{01:02:40--01:04:45}}
  \lecturerresponse{Yes.  Equation~\eqref{eq:l8-generator-rule} defines an
  infinitesimal transformation for every function \(G\), whether or not it is
  conserved.  It becomes a symmetry only when the Hamiltonian is invariant:
  \(\{H,G\}=0\).  The same condition, read in the opposite order, then says
  that \(G\) is conserved.}
\end{classroomqa}

\section{The Angular-Momentum Algebra}

The next calculation lets us stop dragging the particle coordinates through
every rotational problem.  Begin with
\[
  L_x=yp_z-zp_y,\qquad
  L_z=xp_y-yp_x.
\]
In \(\{L_x,L_z\}\), most cross terms vanish because their indices do not form
canonical pairs.  The surviving contractions give
\begin{align}
  \{L_x,L_z\}
  &=\{yp_z-zp_y,xp_y-yp_x\}\\
  &=xp_z-zp_x\\
  &=-L_y.
  \label{eq:l8-lx-lz}
\end{align}
Reversing the order removes the minus sign.  Cycling the coordinates produces
\begin{equation}
  \boxed{\begin{aligned}
    \{L_x,L_y\}&=L_z,\\
    \{L_y,L_z\}&=L_x,\\
    \{L_z,L_x\}&=L_y.
  \end{aligned}}
  \label{eq:l8-angular-algebra}
\end{equation}
In index notation,
\begin{equation}
  \{L_i,L_j\}=\varepsilon_{ijk}L_k.
  \label{eq:l8-angular-algebra-index}
\end{equation}

\begin{figure}[htbp]
  \centering
  \includegraphics[width=0.82\textwidth]{lecture_08_figure_07.png}
  \caption{The cyclic Poisson-bracket algebra of the three angular-momentum
  components.}
  \lectureframe{01:10:52}
\end{figure}

The result could also have been anticipated.  Since \(\boldsymbol L\) is a
vector, bracketing its \(x\)-component with \(L_z\) must give the infinitesimal
change of that component under a \(z\)-rotation, namely \(-L_y\).  The explicit
calculation confirms that geometric argument.  These equations are also the
classical precursors of the quantum angular-momentum commutation relations.

\section{A Fast Gyroscope}

As an application, consider a gyroscope: a rapidly spinning flywheel whose center is a fixed
distance \(R\) from a pivot.  Let \(\boldsymbol r\) point from the pivot to the
wheel center.  The flywheel spins so rapidly about its axle that its dominant
angular momentum remains aligned with that direction.  In the approximation
used in the lecture,
\begin{equation}
  \boldsymbol L=\ell\,\boldsymbol r,
  \qquad
  L_z=\ell z,
  \label{eq:l8-L-parallel-r}
\end{equation}
where \(\ell=L/R\) is treated as constant.  Smaller angular momenta associated
with nutation or motion of the axle are neglected at first.

\begin{figure}[htbp]
  \centering
  \begin{tikzpicture}[scale=0.95,>=Latex]
    \coordinate (P) at (0,3.2);
    \coordinate (C) at (3.7,1.55);
    \draw[->] (0,-0.2) -- (0,4.35) node[above] {\(z\)};
    \fill (P) circle (1.7pt) node[left] {pivot};
    \draw[very thick,->] (P) -- (C)
      node[midway,above right] {\(\boldsymbol r\parallel\boldsymbol L\)};
    \fill (C) circle (1.7pt);
    \draw[dashed] (P) -- (0,1.55);
    \draw[dashed] (0,1.55) -- (C);
    \node[left] at (0,1.55) {\(z\)};
    \begin{scope}[shift={(C)},rotate=-24]
      \draw[very thick] (0,0) ellipse (0.65 and 1.35);
      \draw[thick] (0,0) ellipse (0.42 and 0.9);
    \end{scope}
    \draw[->,thick] (2.15,2.85)
      arc[start angle=140,end angle=35,radius=1.65];
  \end{tikzpicture}
  \caption{Fast-gyroscope approximation.  The pivot-to-center distance is
  fixed, and the dominant spin angular momentum remains aligned with the axle
  and \(\boldsymbol r\).}
\end{figure}

\begin{figure}[htbp]
  \centering
  \includegraphics[width=0.82\textwidth]{lecture_08_figure_08.png}
  \caption{The board sketch of the flywheel together with
  \(H_{\rm rot}=L^2/(2I)\) and the cancellation proving torque-free
  conservation of \(\boldsymbol L\).}
  \lectureframe{01:18:54}
\end{figure}

\subsection{No gravity: why every component stays fixed}

For the simplified axially symmetric rotor,
\begin{equation}
  H_{\mathrm{rot}}
  =\frac{1}{2I}
  \left(L_x^2+L_y^2+L_z^2\right)
  =\frac{L^2}{2I}.
  \label{eq:l8-rotational-hamiltonian}
\end{equation}
The coefficient is \(1/(2I)\), not \(I/2\); that slip is made and corrected
at the board.  Using Eq.~\eqref{eq:l8-angular-algebra},
\begin{align}
  \dot L_x
  &=\{L_x,H_{\mathrm{rot}}\}\\
  &=\frac{1}{2I}\left(
    2L_y\{L_x,L_y\}
    +2L_z\{L_x,L_z\}
  \right)\\
  &=\frac{1}{I}\left(L_yL_z-L_zL_y\right)=0.
  \label{eq:l8-free-lx}
\end{align}
Cyclically,
\begin{equation}
  \dot L_x=\dot L_y=\dot L_z=0.
  \label{eq:l8-free-angular-momentum}
\end{equation}
The scalar \(L^2\) has zero bracket with each component \(L_i\); in the
language of this algebra, it is a Casimir.  Thus a torque-free gyroscope keeps
both the magnitude and direction of its angular momentum.

\begin{classroomqa}
  \audiencequestion{Is the moment of inertia just a number here?
  \lecturetimestamp{01:14:23--01:15:05}}
  \lecturerresponse{For this reduced model, yes.  It encodes the mass
  distribution of the wheel about its spin axis.  A general rigid body has a
  moment-of-inertia tensor and a rotational Hamiltonian
  \(\sum_i L_i^2/(2I_i)\) in principal axes.  The lecture chooses the simpler
  fast, symmetric rotor because the aim is to display the Poisson-bracket
  method rather than develop full rigid-body dynamics.}
\end{classroomqa}

\subsection{Gravity and the sign of the potential}

Take the vertical coordinate \(z\) positive upward and put the pivot at the
origin.  The gravitational potential energy of the wheel center is
\begin{equation}
  V=Mgz.
  \label{eq:l8-mgz}
\end{equation}
Since \(L_z=\ell z\),
\begin{equation}
  V=cL_z,
  \qquad
  c=\frac{Mg}{\ell}=\frac{MgR}{L}.
  \label{eq:l8-potential-lz}
\end{equation}
Here \(c>0\) when \(\boldsymbol L\) is chosen parallel to
\(\boldsymbol r\).  If the flywheel spins in the opposite sense, the signed
relation between \(\boldsymbol L\) and \(\boldsymbol r\) reverses and so does
the corresponding signed precession frequency.

\begin{figure}[htbp]
  \centering
  \includegraphics[width=0.82\textwidth]{lecture_08_figure_09.png}
  \caption{The fixed-length geometry behind
  \(\boldsymbol L=\ell\boldsymbol r\) and \(L_z=\ell z\).}
  \lectureframe{01:25:38}
\end{figure}

\begin{figure}[htbp]
  \centering
  \includegraphics[width=0.82\textwidth]{lecture_08_figure_10.png}
  \caption{An intermediate board state shows \(V=-cL_z\).  The sign is
  corrected later after \(z\) is fixed as upward: with
  \(\boldsymbol L=\ell\boldsymbol r\) and \(\ell>0\), the consistent result is
  \(V=+cL_z\).}
  \lectureframe{01:22:43}
\end{figure}

\begin{classroomqa}
  \audiencequestion{Why may height be replaced by a component of angular
  momentum?  Is the slanted line an axle or a position vector?
  \lecturetimestamp{01:22:20--01:25:43}}
  \lecturerresponse{The slanted line is best labeled
  \(\boldsymbol r\), the vector from the fixed pivot to the wheel center.  A
  physical axle may lie along it, but the relation used in the calculation is
  the fast-spin approximation
  \(\boldsymbol L=\ell\boldsymbol r\).  Their \(z\)-components are therefore
  proportional.  The live discussion changes labels several times before
  settling on this distinction.}
\end{classroomqa}

\begin{classroomqa}
  \audiencequestion{What assumptions make \(\ell\) constant, and what fails if
  the wheel center can slide radially?
  \lecturetimestamp{01:25:43--01:26:38,\allowbreak\ 01:34:22--01:35:00}}
  \lecturerresponse{The dominant spin is assumed much larger than the angular
  momentum associated with the slow motion, and the pivot-to-center distance
  is fixed.  Then the spin magnitude changes little and
  \(\boldsymbol L\) rotates with \(\boldsymbol r\).  If the center could slide
  radially, \(z\) would no longer be proportional to \(L_z\) with one constant
  coefficient, so the reduced Hamiltonian \(V=cL_z\) would not describe the
  motion.}
\end{classroomqa}

\begin{classroomqa}
  \audiencequestion{Why does the board first show a minus sign and later a
  plus sign?  Does the spin direction matter?
  \lecturetimestamp{01:27:13--01:28:15,\allowbreak\ 01:33:09--01:33:21}}
  \lecturerresponse{The initial sign is a live-board error mixed with an
  unsettled axis convention.  Once \(z\) is defined upward,
  \(V=Mgz\) increases with height.  If
  \(\boldsymbol L=\ell\boldsymbol r\) with \(\ell>0\), this gives
  \(V=+cL_z\).  Reversing the spin reverses the signed proportionality and the
  direction of precession, but it does not restore an arbitrary sign choice.}
\end{classroomqa}

\subsection{Precession from three short equations}

The full reduced Hamiltonian is
\begin{equation}
  H=
  \frac{1}{2I}\left(L_x^2+L_y^2+L_z^2\right)+cL_z.
  \label{eq:l8-full-gyro-hamiltonian}
\end{equation}
We have already shown that \(L^2\) has zero bracket with every \(L_i\).  Only
the linear term contributes:
\begin{align}
  \dot L_z&=\{L_z,H\}=0,\\
  \dot L_x&=\{L_x,H\}=c\{L_x,L_z\}=-cL_y,\\
  \dot L_y&=\{L_y,H\}=c\{L_y,L_z\}=cL_x.
  \label{eq:l8-precession-equations}
\end{align}

\begin{figure}[htbp]
  \centering
  \includegraphics[width=0.82\textwidth]{lecture_08_figure_11.png}
  \caption{The final component equations: \(L_z\) is constant while
  \(L_x\) and \(L_y\) rotate into one another.}
  \lectureframe{01:30:35}
\end{figure}

The vertical component and the horizontal magnitude are constant:
\begin{equation}
  \frac{d}{dt}(L_x^2+L_y^2)
  =2L_x(-cL_y)+2L_y(cL_x)=0.
\end{equation}
Combining the horizontal components,
\begin{equation}
  L_+=L_x+iL_y,
  \qquad
  \dot L_+=icL_+,
\end{equation}
so
\begin{equation}
  L_+(t)=L_+(0)e^{ict}.
  \label{eq:l8-precession-solution}
\end{equation}
The angular-momentum vector therefore precesses around the vertical axis with
angular speed \(|c|=MgR/L\).  Its tilt remains fixed in this idealized steady
motion.

\begin{classroomqa}
  \audiencequestion{Was the rotational-energy term neglected when the
  precession equations were written?
  \lecturetimestamp{01:35:01--01:35:49}}
  \lecturerresponse{No.  It belongs in the full Hamiltonian as in
  Eq.~\eqref{eq:l8-full-gyro-hamiltonian}.  It simply contributes zero to
  \(\dot L_i\) because \(\{L_i,L^2\}=0\), a cancellation already proved in the
  torque-free calculation.  The student question makes this implicit step
  explicit.}
\end{classroomqa}

\begin{classroomqa}
  \audiencequestion{Is uniform precession an equilibrium?  What happens after
  a small kick? \lecturetimestamp{01:32:38--01:33:08}}
  \lecturerresponse{It is better called a steady or stationary motion than a
  static equilibrium.  A small kick generally adds nutation, an oscillation of
  the tilt angle.  For a very rapidly spinning wheel that oscillation is small,
  which is why the reduced treatment keeps only nearly uniform precession.}
\end{classroomqa}

\begin{classroomqa}
  \audiencequestion{Does precession require a support or a real axle?
  \lecturetimestamp{01:31:49--01:32:38,\allowbreak\ 01:33:49--01:35:00}}
  \lecturerresponse{This mechanical model assumes a fixed pivot-to-center
  distance, so some constraint supplies the support.  The deeper algebraic
  result is broader: whenever an angular momentum has a Hamiltonian term
  linear in one component, the same bracket algebra produces precession.
  The microscopic source of the angular momentum need not be a literal wheel
  on an axle.}
\end{classroomqa}

\begin{classroomqa}
  \audiencequestion{Is this the same precession as electron or nuclear spin in
  a magnetic field, or as a planet's orbit precessing?
  \lecturetimestamp{01:35:50--01:37:00}}
  \lecturerresponse{A magnetic moment coupled to a magnetic field has a
  Hamiltonian linear in a spin component and therefore obeys the same
  angular-momentum algebra.  Nuclear-spin precession underlies nuclear magnetic
  resonance.  The precession of a planetary orbit's perihelion is a different
  phenomenon that happens to share the word \emph{precession}; it is not the
  same calculation.}
\end{classroomqa}

\begin{classroomqa}
  \audiencequestion{Is the side-mounted gyroscope essentially the same object
  as a spinning top? \lecturetimestamp{01:37:02--01:37:23}}
  \lecturerresponse{Yes.  The distinction is mostly geometric language: when
  the spinning axis extends sideways from its support one usually says
  gyroscope; when it is nearly vertical one says top.  The angular-momentum
  and torque structure is the same.}
\end{classroomqa}

\section{What the Calculation Accomplished}

No Euler angles were introduced, and no coordinates for every piece of the
wheel were followed.  Once
\[
  \{L_i,L_j\}=\varepsilon_{ijk}L_k
\]
was known, the equations for the variables of interest followed directly.
That is the intended demonstration of power.  The oscillator showed that
Poisson brackets can reproduce familiar dynamics.  Angular momentum showed
that they do much more:
\begin{itemize}
  \item a conserved quantity generates its associated infinitesimal symmetry;
  \item conservation of \(G\) and invariance of \(H\) are the two readings of
  \(\{G,H\}=0\);
  \item the algebra of observables can close on itself, allowing dynamics to be
  solved without returning to microscopic coordinates;
  \item a Hamiltonian linear in one angular-momentum component produces
  precession around that axis.
\end{itemize}
The lecture stops before the full Euler equations for a general rigid body.
Its narrower conclusion is already complete: the right generators and their
Poisson algebra expose both symmetry and motion with remarkable economy.
